%% ====================================================================
\section*{Estructura conceptual del TFM}
\addcontentsline{toc}{section}{Estructura conceptual del TFM}
\label{sec:estructura-conceptual}
%% ====================================================================

Este TFM estudia un problema de coordinación distribuida en almacenes industriales:
varios AGVs deben decidir, sin un asignador central permanente, qué cargas atender, con
qué coalición, por qué ruta y bajo qué restricciones físicas. La dificultad no está solo
en escoger una tarea. Una carga pesada o voluminosa puede requerir un quorum mínimo de
robots, una distribución espacial compatible con el contacto, suficiente batería para
volver a una zona segura, límites de torque en rueda, evitación de obstáculos y
capacidad de mantener la formación durante el transporte.

\subsection*{Problema de coordinación}
\addcontentsline{toc}{subsection}{Problema de coordinación}

El sistema se compone de $N$ robots móviles y $K$ tareas activas. Cada tarea $k$ tiene
una posición, una prioridad, una demanda de servicio y una carga física asociada. En el
modelo homogéneo, la demanda se resume en un número mínimo $n_k$ de robots. En el modelo
físico, la demanda se escribe como una capacidad requerida $D_k$ y la coalición aporta
una capacidad efectiva $S_k$ que depende de la masa, la geometría, la fricción, el
contacto, la energía disponible y la orientación de la carga. La condición de servicio
deja de ser únicamente $|\mathcal C_k|\ge n_k$ y pasa a ser
\[
  S_k(\mathcal C_k,t)\ge D_k(t).
\]
Esta desigualdad es el núcleo del problema: una coalición no es válida porque tenga
muchos robots, sino porque cierra las restricciones físicas necesarias para mover la
carga con seguridad.

\subsection*{Teoría de juegos}
\addcontentsline{toc}{subsection}{Teoría de juegos}

La teoría de juegos ofrece el lenguaje para describir decisiones acopladas entre
agentes. Cada robot es un jugador, cada tarea es una estrategia y el valor de elegir una
tarea se mide con un \emph{payoff}. A diferencia de una asignación centralizada, aquí el
payoff de un robot no depende solo de la tarea, sino también de cuántos robots ya la
atienden, de la distancia, de la conectividad local, de la batería y de la contribución
mecánica que ese robot puede aportar.

El punto importante es que el juego no se usa como metáfora económica, sino como una
estructura de control. Si el payoff de una tarea aumenta cuando falta capacidad y
disminuye cuando la tarea queda cubierta, entonces la flota recibe una señal distribuida
de reclutamiento. Los robots cercanos a una carga deficitaria tienden a migrar hacia
ella; los robots que ya no aportan valor marginal tienden a liberar capacidad para otras
tareas. En el caso integrable, estos payoffs derivan de un potencial global, lo que
permite demostrar propiedades de estabilidad y convergencia.

\subsection*{Dinámicas poblacionales}
\addcontentsline{toc}{subsection}{Dinámicas poblacionales}

Las dinámicas poblacionales convierten el juego en una ecuación de movimiento sobre el
símplex de decisiones. Cada robot mantiene un vector
\[
  y_i=(y_{i0},y_{i1},\ldots,y_{iK})\in\Delta^{K+1},
\]
donde $y_{ik}$ representa su preferencia por la tarea $k$ y $y_{i0}$ representa
inactividad o reserva. La dinámica de Smith transfiere masa decisional desde estrategias
con menor payoff hacia estrategias con mayor payoff mediante comparaciones por pares:
\[
  \dot y_{ik}
  =
  \sum_{\ell=0}^{K} y_{i\ell}[F_{ik}-F_{i\ell}]_+
  -
  y_{ik}\sum_{\ell=0}^{K}[F_{i\ell}-F_{ik}]_+.
\]
Esta ecuación preserva el símplex, evita probabilidades negativas y permite que cada
robot revise su decisión con información local. En la implementación, la preferencia
continua se proyecta a una decisión ejecutable mediante compromiso, histéresis, guardias
de quorum y \emph{clearing} entero.

\subsection*{MARL como referencia algorítmica}
\addcontentsline{toc}{subsection}{MARL como referencia algorítmica}

El aprendizaje por refuerzo multiagente (MARL) plantea el mismo problema desde otra
perspectiva. En lugar de diseñar explícitamente el payoff y demostrar la dinámica, cada
robot aprende una política $\pi_i(a_i\mid o_i)$ que asigna acciones a observaciones para
maximizar una recompensa acumulada. MARL es atractivo porque puede capturar patrones
difíciles de modelar a mano, pero también introduce dependencia del entrenamiento,
dificultad de generalización y menor interpretabilidad formal.

Por ese motivo, MARL se usa en este TFM como referencia comparativa y como inspiración
para métricas, no como sustituto del modelo analítico. La propuesta Smith-QR conserva
ecuaciones trazables: si un robot cambia de tarea, se puede explicar por déficit de
capacidad, precio temporal, distancia, conectividad, batería o restricción de seguridad.
Esta trazabilidad es clave en robótica industrial, donde no basta con que una política
funcione en promedio; también debe explicar por qué una coalición se forma y bajo qué
condiciones deja de ser segura.

\subsection*{Aporte integrado}
\addcontentsline{toc}{subsection}{Aporte integrado}

La contribución del TFM es unificar seis capas que suelen estudiarse por separado:
\[
  \begin{aligned}
  \text{juego}
  &\rightarrow
  \text{dinámica poblacional}
  \rightarrow
  \text{capacidad mecánica}\\
  &\rightarrow
  \text{control vectorial ejecutable}
  \rightarrow
  \text{escalabilidad}
  \rightarrow
  \text{validación multi-métrica}.
  \end{aligned}
\]
La unión entre capas se realiza mediante una cadena de residuos: déficit de capacidad,
déficit de wrench, balance port-Hamiltoniano, margen de batería y congestión prevista.
Así, una capa no queda justificada por su elegancia aislada, sino por la señal que entrega
a la siguiente y por el gate que puede falsarla.
La primera capa decide qué tareas son atractivas. La segunda transforma payoffs en
revisión distribuida. La tercera reemplaza el conteo simple de robots por capacidad
efectiva, torques, contacto, formación y batería. La cuarta baja esas variables a un
bucle finito con contornos circulares/superelípticos, firma fuerza--torque, wrench,
fases, rigidez, CBF y uniciclo. La quinta separa
dos lecturas: aproximación determinista de Smith para conectar la ODE con flotas finitas,
y MFG como arquitectura futura para campos de densidad y peajes de congestión. La sexta
evalúa no solo entregas y distancia, sino también energía, vida útil de batería,
saturación de actuadores, seguridad, congestión y robustez ante fallos.

Con esta organización, Smith-QR se interpreta como el núcleo ejecutable de un framework
más amplio: una arquitectura de coordinación distribuida donde la teoría de juegos
define incentivos, las dinámicas poblacionales producen adaptación continua, la mecánica
garantiza que las coaliciones sean físicamente útiles, el control vectorial traduce esas
coaliciones a comandos ejecutables y la aproximación determinista explica cuándo la ODE
representa a la flota finita. Las ideas MFG quedan como ruta de escalado por campos
escalares, sin convertirlas en garantía experimental central del TFM. Los gates G3 y G4
validan las conexiones mínimas del motor avanzado: capacidad wrench frente a cardinalidad,
firma fuerza--torque, peajes predictivos, batería y balance energético.
